\clearpage
\SourceLeaf{377}%
\begin{center}\large [CONFORMAL MAPPING. MEASURE OF CURVATURE.]\end{center}
\begin{center}[1.]\end{center}
\begin{center}\textsc{Gauss} to \textsc{Schumacher}. Göttingen, 5 July 1816.\end{center}

\mbox{. . . . . . .} I have not yet seen the announcement of your society’s prize question. I also discussed with \textsc{Lindenau} a prize question that is to be set in the new journal with a prize of 100 ducats. An interesting problem had occurred to me, namely
\begin{quote}
\noindent “\textbf{in general}, to project (map) a given surface onto another\\
“(given) surface so that the image is similar to the original\\
“in its smallest parts.”
\end{quote}

A special case occurs when the first surface is a sphere and the second is a plane. Here the stereographic and Mercator projections are particular solutions. What is wanted, however, is the \textls[120]{general} solution that includes all the particular ones, for every kind of surface.

\textsc{Monge} and \textsc{Ponsot} are said to have already worked on this in the Journal philomathique (as \textsc{Burckhardt} wrote to \textsc{Lindenau}). Since I do not know exactly where, I have not yet been able to look it up and therefore do not know whether those gentlemen’s solutions fully agree with my idea and exhaust the subject. If they do not, this would seem to me a suitable prize question for a society on some occasion. At the society here, my turn to set a question comes only once every 12 years. \mbox{. . . . . . .}
